\section{Method}
\label{sec:method}

\subsection{Frozen-backbone VLM with a trainable bridge}
\label{sec:method-arch}

Our base model follows Vintern-1B: a vision encoder \textbf{InternViT-300M} and a
language model \textbf{Qwen2-0.5B}, connected by a projector. Unlike Vintern-1B's
own downstream-adaptation recipe (full fine-tuning of the encoder and projector,
LoRA on the decoder), we freeze both backbones and replace only the projector
with a trainable \emph{bridge} module. Training optimises the bridge alone, by
cross-entropy on a reference answer (Fig.~\ref{fig:method}).

\begin{figure}[t]
  \centering
  \includegraphics[width=\linewidth]{fig_method}
  \caption{The frozen-backbone architecture. An image is encoded as $1$--$T$
  tiles by a frozen InternViT-300M; a trainable bridge (green, ${\sim}1\%$ of
  parameters) maps the patch tokens to $k$ vision tokens; a frozen Qwen2-0.5B
  decoder --- optionally with a rank-16 attention LoRA --- generates the
  Vietnamese answer.}
  \label{fig:method}
\end{figure}

We study five bridges spanning a capacity ladder from a one-token linear
projector to a 16-query Q-Former (Table~\ref{tab:bridges}). The
\textbf{multi-token bridge} --- eight mean-pooled tokens, one anchor plus seven
semantic --- is our default: it has the lowest trainable-parameter count among
the multi-token designs (0.78\,\%) and the lowest validation cross-entropy of all
five.

\begin{table}[t]
  \centering
  \caption{The five bridge designs, ordered by trainable-parameter count.}
  \label{tab:bridges}
  \begin{tabular}{lccl}
    \toprule
    Bridge & $k$ tokens & Trainable params & Mechanism \\
    \midrule
    Residual        & 1  & 4.86\,M (0.52\,\%) & linear projector $+$ LayerNorm/GELU residual \\
    \textbf{Multi-Token} & 8  & 7.35\,M (0.78\,\%) & mean-pool patches $\rightarrow$ 8 tokens (1 anchor $+$ 7 semantic) \\
    Tile-Attention  & 8  & 4.14\,M (0.44\,\%) & dense patch self-attention, then pool \\
    Light Q-Former  & 8  & 27.6\,M (2.87\,\%) & 8 learned queries, 2-layer cross-attention \\
    Full Q-Former   & 16 & 69.4\,M (6.91\,\%) & 16 learned queries, 4 layers, image--text fusion \\
    \bottomrule
  \end{tabular}
\end{table}

\subsection{Two knobs, one diagnostic ladder}
\label{sec:method-ladder}

Given a frozen backbone, a small additional parameter or compute budget can be
spent on the \textbf{vision side} or the \textbf{language side} of the pipeline.
We work through six research questions, one lever at a time, each measured as
\dF{} against a fixed anchor (the multi-token bridge, one tile, no LoRA):

\begin{table}[t]
  \centering
  \caption{The six research questions. RQ1--RQ4 are vision/training-side; RQ5
  straddles; RQ6 is the language side.}
  \label{tab:rqs}
  \begin{tabular}{clc}
    \toprule
    & Question & Lever \\
    \midrule
    RQ1 & Does a larger / more expressive bridge close the token-F1 gap? & bridge capacity \\
    RQ2 & Is the bridge itself the constraint, or already near-optimal? & validation CE across bridges \\
    RQ3 & Does more visual resolution --- more image tiles --- help? & \ntiles{} \\
    RQ4 & Does \emph{adaptive}, per-question visual-compute allocation help, driven by question type? & router $+$ policy \\
    RQ5 & Does a richer training signal or explicit representation alignment help? & sampling, projector-KD \\
    RQ6 & Does adding capacity to the \emph{frozen decoder} help, and where? & attention vs.\ FFN LoRA \\
    \bottomrule
  \end{tabular}
\end{table}

The shape of the answers --- which levers move F1 and which do not --- is the
paper's main finding (Sec.~\ref{sec:abl-summary}).

\subsection{Adaptive visual computation (RQ3--RQ4)}
\label{sec:method-adaptive}

\paragraph{Tile lever.}
An image is split into $\ntiles \in \{1, 3, 6\}$ tiles of $448{\times}448$; each
tile is a separate InternViT forward pass, so \ntiles{} is the primary
visual-compute knob (Sec.~\ref{sec:abl-compute} characterises its FLOPs/latency
cost).

\paragraph{Router.}
Running in parallel with --- and far more cheaply than --- the vision encoder,
the router produces two signals. A \textbf{cognitive prior} $P(r \mid Q)$: a
PhoBERT head over AutoViVQA's eight reasoning types, seeing the question text only
(validation macro-F1 $0.91$; since the reasoning-type label is largely derivable
from question surface form, this is effectively a question-pattern prior). A
\textbf{cheap visual state} $f(I, Q)$: the InternViT CLS embedding at one tile
(PCA to 64 dims) plus question length and three image scalars (clarity,
occlusion, object density).

\paragraph{Offline oracle-guided policy.}
For a cost trade-off $\lambda$, define the oracle action over the discrete space
$a = (\ntiles, \text{bridge})$ as
$a^\ast(x, \lambda) = \arg\max_a \left[ M(a; x) - \lambda \cdot C(a) \right]$,
where $M(a; x)$ is per-sample CIDEr against the five references and
$C(a) = \ntiles / \max\text{tiles}$. We evaluate \emph{all} actions on
\emph{every} train/val/test question once (the \emph{oracle sweep},
Sec.~\ref{sec:setup-hw}), then train a policy MLP
$\pi_\theta(P(r{\mid}Q), f(I,Q), \lambda)$ by supervised classification against
$a^\ast$ on the training split. Ablation arms differ only in the policy's inputs:
reasoning-type only, visual-state only, both, and fixed / random baselines. RQ4
reduces to whether any learned policy beats the fixed policy ``best bridge at one
tile'' on held-out test.

\subsection{Training-signal and alignment interventions (RQ5)}
\label{sec:method-rq5}

Two ways to give the bridge a better target without touching the backbones:
\textbf{(a) multi-reference answer sampling} --- instead of always training toward
the first of the five references, resample a reference each epoch;
\textbf{(b) projector-level knowledge distillation} --- add a KD term pulling the
bridge output toward Vintern-1B's own pre-aligned \texttt{mlp1} projector, at the
feature level (\texttt{align-feat}) or the decoder-logit level
(\texttt{align-logit}), with weight $\alpha$ on the KD term.

\subsection{Decoder LoRA as a targeted intervention (RQ6)}
\label{sec:method-rq6}

Finally, the smallest departure from ``backbone frozen'': a rank-16 LoRA
($\alpha = 32$, 2.16\,M parameters, $0.23\%$ of the model) on the frozen Qwen2
decoder, trained on top of a fixed bridge. We vary \textbf{which modules} the
LoRA touches --- the attention projections $\{q,k,v,o\}$, the feed-forward
projections $\{\text{gate},\text{up},\text{down}\}$, or all seven --- to localise
where the decoder's useful headroom lies, and vary the number of LoRA epochs
($1$, $3$) to bound the effect.
